如果两个方案都可行，把它们按比例混合后是否仍可行？凸集合就是对这个问题给出肯定答案的集合。这个看似简单的几何性质，保证了朝更优方案走一小步不会突然越出可行域，是``局部最优等于全局最优''证明中不可缺少的一半（另一半由凸目标函数承担，在下一单元展开）。

本单元沿三层几何逐步推进：

\begin{enumerate}
\tightlist
\item
  线段、凸组合与凸包从两点混合出发，建立凸集的定义、构造方法与反例判断。
\item
  等式与不等式如何塑造可行域解释仿射集、超平面、半空间和多面体，并用分离超平面说明外部点为何能被线性证据排除。
\item
  从方向的集合到半正定锥把非负数推广到锥与矩阵，建立二阶锥、半正定锥和线性矩阵不等式的直觉。
\end{enumerate}

读完后，面对一个约束集合，应能优先尝试``由已知凸集通过交、仿射像或原像构造''的证明；若怀疑它非凸，则寻找两点及其中间点作为反例。

\subsection{一个直接的凸性检验}

给定平面上三个约束：\(x\ge0,\;y\ge0,\;x+y\le10\)。任取两个可行点，如 \((2,3)\) 和 \((7,1)\)，\(\theta\) 取 \(0.4\)，混合点 \((0.4\times2+0.6\times7,\;0.4\times3+0.6\times1)=(5,1.8)\)。逐个验证：\(5\ge0\)，\(1.8\ge0\)，\(5+1.8=6.8\le10\)，均成立。这三个半空间的交集是凸三角形，这是判断可行域凸性的典型图景。若加入 \(x^2+y^2=5\) 作为新约束，两个可行点 \((1,2)\) 和 \((2,1)\) 的中点是 \((1.5,1.5)\)，其坐标平方和 \(4.5\neq5\)，破坏了凸性：圆含在约束中使可行域变成两段弧线，它们的凸组合留下弧形缺口，这正是等式使得可行域不凸的直观解释。

\subsection{怎么读这一章}

先不要把凸集理解成某种特定外形。读第一篇时只问一件事：两个可行方案按比例混合后，是否仍然可行。第二篇再看线性等式和不等式怎样制造最常见的凸集。第三篇把同一个``非负组合仍可行''的想法推广到方向和矩阵。若还不清楚这项几何性质为何重要，可先回看 \hyperref[ch:071]{``从现实选择到优化问题''}中局部最优与全局最优的区别。

\subsection{本章篇目}

以下篇目按概念依赖排列。

\begin{itemize}
\tightlist
\item \hyperref[sec:09-Convex-Optimization/02-Convex-Sets/01]{``线段、凸组合与凸包''}；
\item \hyperref[sec:09-Convex-Optimization/02-Convex-Sets/02]{``等式与不等式如何塑造可行域''}；
\item \hyperref[sec:09-Convex-Optimization/02-Convex-Sets/03]{``从方向的集合到半正定锥''}。
\end{itemize}
